\chapter{Introduction}
\label{introduction.ch}

High-performance computing (HPC) environments, such as those at the Ohio Supercomputer Center (OSC) and the Texas Advanced Computing Center (TACC), are essential for running computationally intensive tasks like weather simulations, genome sequencing, and deep neural network (DNN) training. These shared-resource systems allow multiple users to run jobs with varying resource demands and execution times. While HPC platforms offer different interfaces for resource requests, efficient utilization depends on users' ability to assess workload characteristics and request appropriate resources accurately.
However, managing jobs in HPC environments presents several challenges, often leading to inefficient resource allocation. Users must navigate complex configurations, adapt to evolving infrastructure, and work with limited estimation tools, frequently requiring manual adjustments. Continuous monitoring and repeated job resubmissions add to inefficiencies, extending execution times. These challenges include:

\begin{itemize}
    \item Resource Complexity: Researchers need a deep understanding of available resources, allocation policies, limitations, costs, and configuration options to utilize HPC resources effectively.
    \item Inconsistent Job Configurations: Simple estimations become unreliable due to changes in code, parameter settings, software, or hardware upgrades.
    \item Steep Learning Curve: Navigating profiling and analytics tools to understand job requirements and optimize resource allocations requires significant training and effort.
    \item Constant Updates: Researchers must stay informed about cyberinfrastructure (CI) and middleware changes to maintain optimal resource utilization.
    \item Limited Resource Estimators: Current tools do not generate execution plans that account for available architectures or CI batch policies, necessitating manual adjustments.
    \item Continuous Monitoring: Researchers must constantly monitor progress and performance despite pre-profiling jobs, often needing to resubmit jobs to improve efficiency or extend wall time.
    \item Human Intervention: The frequent need to separate jobs into smaller jobs to fit CI allocation policies leads to a longer time to completion due to the need for human monitoring and intervention.
\end{itemize}

These challenges often result in researchers over-provisioning resources to ensure job completion, leading to longer queue wait times, inefficient resource use, and increased costs. Chapter \ref{problem.ch} describes the motivation for creating AI-enabled frameworks to address these challenges. 

This dissertation explores and presents the integration of machine learning into HPC cyberinfrastructure to enhance resource provisioning. AI-driven solutions can help users better understand their resource requirements, enabling more informed allocation decisions and improving overall efficiency. We aim to address resource management challenges in HPC environments by developing intuitive and programmable AI-based methods.


Our research introduces a comprehensive framework that utilizes two key strategies:
\begin{itemize}
    \item Application-Specific Resource Estimation: We developed the HPC Application Resource Predictor (HARP). This advanced tool leverages machine learning to create tailored models for estimating the resource requirements of individual applications. HARP addresses the complexity of modern HPC applications by considering diverse factors such as application configurations, input size, and hardware variability. This approach significantly improves traditional estimation methods by adapting models to each application's unique characteristics. The reference architecture and the components used to create HARP are described in Chapter \ref{OverviewSolution.ch} Section \ref{HARP}.
    \item Intelligent CI-Aware Scheduling: This strategy produces application-specific estimates made specific to  Cyberinfrastructure (CI) configurations. Doing so enables context-aware execution recommendations across diverse HPC systems that are feasible based on the CI budget limitations and costing policies. This intelligent scheduling system considers the estimated resource requirements and the current state of the HPC environment in terms of available resources to schedule or reschedule jobs. The reference architecture and the components used to create the Intelligent Scheduler are described in Chapter \ref{OverviewSolution.ch} Section \ref{iSchedulerF}.
\end{itemize}
    

Key innovations of our research include:
\begin{itemize}
    \item Cyberibnfrastructre Configuration Database: We constructed a comprehensive database of HPC system configurations by combining manually scraped CI documentation with programmatically generated information. This database is crucial in assessing resource suitability and generating viable execution plans. We call this database a pre-SLURM validation database. Chapter \ref{OverviewSolution.ch} explains the development of a database that supplies data to the estimation engine, assists in generating execution plans, and evaluates the feasibility of executions.
    \item Efficient Training Data Generation: We developed a novel method to generate application-specific training data through downsized execution campaigns. This approach reduces the cost and time required for data collection by a factor of 7, making it feasible to create accurate models for a wide range of applications. Chapter \ref{GenerateData.ch} describes how we emulate application profiling to generate scalable training datasets for resource estimations. 
    \item DNN-Estimator for AI Workloads: Recognizing the growing importance of deep learning in scientific research, we created a specialized estimator for deep neural network (DNN) training and inferencing tasks. The DNN-Estimator can predict resource requirements based on network architecture, dataset characteristics, and hardware configurations. Chapter \ref{FutureWork.ch} describes the white box-model aware resource estimators for training the model. 
    \item Intelligence Plane: This component dynamically generates execution plans using the pre-SLURM database and resource estimators. Integrated with TAPIS, it can schedule, monitor, and reschedule jobs in real-time, adapting to changing system conditions and job performance. The reference architecture and the components used to create the Intelligent Scheduler are described in Chapter \ref{OverviewSolution.ch}.
\end{itemize}


To demonstrate the real-world applicability of our framework, we present a detailed use case in animal ecology primarily focused on workflow interactions with centers like OSC or TACC. This pipeline begins with data collection field sensors or cameras and is analyzed by AI inference models deployed on edge, triggering a chain of automated processes for model retraining for better accuracy. Once new data is transferred to the center, the Intelligence Plane (IP) orchestrates tasks, such as data labeling and model retraining. The IP employs sophisticated resource estimation techniques, using the DNN Estimator or HARP, to determine optimal resource allocation. It then validates these requirements against a pre-SLURM database and generates execution plans accordingly. The IP continuously monitors job progress, dynamically rescheduling tasks to ensure efficient resource utilization. 

Our future work will proceed in two directions: 1. enhancing DNN estimator efficiency, and 2. improving user interaction through chatbot integration. Our DNN execution time estimation framework currently relies on generating High-Level Optimization (HLO) graphs for a given DNN model architecture (DNN model summary)—an intermediate representation used by compilers like XLA in machine learning frameworks—to predict training times. While HLO graphs effectively optimize computations for different hardware accelerators, our current approach faces a significant bottleneck: we must execute each new DNN model on target hardware platforms to generate these graphs before estimating training times. For example, estimating VGG16 training time on a V100 GPU requires running the model for one step on that specific hardware to generate the HLO graph, from which we extract features for our time estimation model. This dependency necessitates maintaining short-lived queues across various systems just for graph generation. To overcome this limitation, we propose using generative AI models to dynamically create HLO graphs for given DNN architectures and hardware configurations, eliminating the need for actual hardware execution and significantly reducing computational overhead. Additionally, we intend to develop an interactive chatbot by integrating our iScheduler helper framework with an Ollama model. This chatbot will enable users to query models, explore configurations, and receive real-time feedback on resource selection and feasibility options. Essentially, we aim to streamline the resource estimation process and provide a more intuitive, chat-based assistant for resource management in HPC environments, thereby enhancing the user experience while improving system efficiency.


